% Figure: Saint-Venant's principle as a decay of the non-uniform part of
% the stress with distance from a concentrated load. A bar is loaded at
% one end by a point load; three cross-sections are sketched with their
% stress profiles: strongly peaked at the load, less peaked one width
% away, essentially uniform beyond one width. The right panel plots the
% peak-to-mean stress ratio against distance in bar-widths, showing the
% ratio decaying toward 1 (uniform) within about one width. The decay is
% drawn as an exponential toward the uniform value, the standard
% engineering picture of the principle.
\begin{figure}[htbp]
\centering
\begin{tikzpicture}
% left panel: bar with point load and three section profiles
\begin{scope}[shift={(0,0)}]
% bar outline
\draw[engblack, thick] (0,-1) rectangle (5,1);
% point load at left end
\draw[engred, very thick, -{Latex}] (-1.2,0) -- (0,0);
\node[engred, font=\scriptsize, anchor=south] at (-0.6,0.05) {$P$};
% three section lines
\foreach \x/\lab in {0.6/A, 2.0/B, 4.2/C} {
  \draw[enggray, dashed] (\x,-1) -- (\x,1);
  \node[enggray, font=\scriptsize, anchor=south] at (\x,1.02) {\lab};
}
% profile at A: strongly peaked (near the load)
\draw[engblue, very thick] plot[domain=-1:1,samples=30] ({0.6+0.9*exp(-6*\x*\x)},{\x});
% profile at B: milder peak
\draw[engblue, very thick] plot[domain=-1:1,samples=30] ({2.0+0.5*(0.55+0.45*exp(-2*\x*\x))},{\x});
% profile at C: nearly uniform
\draw[engblue, very thick] (4.2+0.42,-1) -- (4.2+0.42,1);
\node[engblack, font=\scriptsize, anchor=north] at (2.5,-1.15) {one bar-width $\approx$ B};
\end{scope}
% right panel: decay curve
\begin{scope}[shift={(7.5,0)}]
\begin{axis}[
  width=6.2cm, height=6.0cm,
  xlabel={distance from load, $x/w$},
  ylabel={peak/mean stress},
  xmin=0, xmax=3, ymin=0.8, ymax=3.2,
  axis lines=left,
  tick label style={font=\scriptsize},
  label style={font=\scriptsize},
  at={(0,-3cm)},
]
% peak/mean ~ 1 + 2 exp(-2.5 x): 3 at x=0, ~1.16 at x=1, ~1 beyond
\addplot[engblue, very thick, domain=0:3, samples=60] {1 + 2*exp(-2.5*x)};
\addplot[enggray, dashed, samples=2, domain=0:3] {1};
\node[enggray, font=\scriptsize, anchor=south east] at (axis cs:3,1) {uniform};
\filldraw[engorange] (axis cs:1,1.164) circle[radius=1.6pt];
\node[engorange, font=\scriptsize, anchor=south west] at (axis cs:1.05,1.16) {$x=w$};
\end{axis}
\end{scope}
\end{tikzpicture}
\caption[Saint-Venant's principle as stress-uniformity recovery]{The
stress profile across a bar loaded by a concentrated force at one end.
Close to the load (section A) the profile is sharply peaked and depends on
exactly how the load is applied; one bar-width away (section B) the peak
has largely smoothed; beyond about one width (section C) the profile is
uniform and the cross-section knows only the load's resultant. The right
panel plots the ratio of peak stress to mean stress against distance in
bar-widths: the non-uniform part decays quickly, and by one width the
overshoot is already small. This decay is what licenses the uniform
$\sigma = N/A$ away from load points and confines the awkward local
analysis to a boundary layer one width deep.}
\label{fig:vol03ch03:saint-venant-decay}
\end{figure}
